\section{What The Horizon Carries Forward}

This chapter has forced surface valuation. It began with alpha as a finite
receipt inherited from Chapter 09. That receipt could travel because it
carried its seed, residue, rung, extrapolation bound, and path cost. It
could not become a hidden interior possession. The horizon made that
discipline unavoidable.

The boundary radiation residue was first pushed toward a lightlike
approach. The relativistic orbit coefficient read the cost of preserving the
invariant as the surface approached the null face. The event horizon study
certificate then bounded that approach by finite probes and carried their
cost. The grammar permitted timelike approach and forbade a completed
crossing.

The Cooper-pair horizon then masked the interior. An ordinary spin-balanced
Cooper-pair interior and a neutrino-exclusion interior remained different
as hidden stories, but the exterior could not spend that difference once
their exported surfaces agreed. The horizon quotient identified interiors
only by exact equality of surface observables.

The surface tuple was narrow: mass, charge, spin-gravity residue, and alpha.
Mass carried second-difference strain. Charge carried loop count.
Spin-gravity residue carried the cost-bearing correction of horizon
approach and exclusion drift. Alpha carried the bounded finite receipt. The
pairing witness, spin witness, exclusion witness, and interior path history
were masked.

The neutrino then appeared as a replicating portfolio. It was not an opened
interior. It was the exclusion label carried by exact replication of the
exported surface payoffs. If mass, charge, spin-gravity residue, and alpha
match, the portfolio stands in for the unobservable interior. If any surface
observable fails to match, the portfolio fails.

Color controlled leakage. A color residue at the single surface correction
would have meant that the horizon had imported an interior distinction under
a placement label. The grammar therefore forced color to be null. Flavor
then carried the allowed correction: the difference between reader-seeded
horizon alpha and charge-only alpha. Subtracting that residue restored the
charge-only receipt without opening the horizon.

Finally, the QFT alpha receipt scaled the single flavor residue across the
finite sector certificate. Three color faces, six flavor labels, and two
phase labels gave thirty-six sectors. Color counted without leaking, flavor
supplied the residue, and phase supplied the orientation faces. The receipt
remained a surface receipt, not a completed interior field.

What passes forward is therefore a horizon-valued alpha. The next chapter
does not inherit a story about what was hidden behind the surface. It
inherits the surface receipt itself: mass, charge, spin-gravity residue,
alpha, the color null control, the flavor exclusion residue, and the
thirty-six-sector scaling. The grammar has learned how to value what it
cannot inspect.
